%%
%% winter-lecture25.tex
%% 
%% Made by Alex Nelson <pqnelson@gmail.com>
%% Login   <alex@lisp>
%% 
%% Started on  2026-03-03T09:50:32-0800
%% Last update 2026-03-03T09:50:32-0800
%% 

\lecture{}

\begin{node}
Compact operators are the norm closure of finite-rank oeprators. What
do they look like in the wild?
\end{node}

\begin{definition}
Let $\mathcal{H}=L^{2}(\RR^{d})$.
\begin{enumerate}
\item An \define{Integral Operator} is given by
  \begin{equation*}
(Tf)(x)=\int_{\RR^{d}}K(x,y)f(y)\,\D y
  \end{equation*}
  where $K(x,y)$ is its \define{Kernel}.
\item If $K\in L^{2}(\RR^{d}\times\RR^{d})$ is the kernel of an
  integral operator $T$, then we say $T$ is a \define{Hilbert--Schmidt Operator}.
\end{enumerate}
\end{definition}

\begin{fact}
Integral operators on Hilbert spaces are linear.
\end{fact}

\begin{proposition}
Hilbert--Schmidt operators are compact.
\end{proposition}

\begin{proof}
We will prove it is bounded in order to prove it is compact. Remember,
bounded is equivalent to continuous (in this setting).

\textsc{Claim 1:} $T$ is bounded. By Fubini, $y\mapsto|K(x,y)|^{2}\in L^{1}$
for almost every $x$ since
\begin{equation}
\iint_{\RR^{d}\times\RR^{d}}|K|^{2}=\int_{\RR^{d}}\underbrace{\left(\int_{\RR^{d}}|K|^{2}\,\D y\right)}_{\text{finite for a.e. }x}\D x
\end{equation}
Then by Cauchy-Schwarz,
\begin{equation}
|\int K(x,y)f(y)\,\D y|\leq(\int|K|^{2}\,\D
y)^{1/2}(\int|f(y)|^{2}\,\D y)^{1/2},
\end{equation}
then squaring both sides gives us
\begin{equation}
\|Tf\|_{L^{2}}^{2}
\leq\int[\int|K(x,y)|^{2}\,\D y\int|f(y)|^{2}\,\D y]\D x
=\|K\|_{L^{2}(\RR^{d}\times\RR^{d})}^{2}\|f\|_{L^{2}(\RR^{d})}^{2}.
\end{equation}
Then
\begin{equation}\label{eq:math-108b:winter2026:lecture25:star}
\|T\|_{L^{2}}\leq\|K\|^{2}.
\end{equation}
Hence the operator is bounded.

\textsc{Claim 2:} $T$ is compact. Let $\{\varphi_{k}\}_{k\in\NN}$ be
an orthonormal basis for $L^{2}(\RR^{d})$. Then
$\{\varphi_{k}(x)\varphi_{\ell}(y)\}_{k,\ell\in\NN}$ is an orthonormal
basis for $L^{2}(\RR^{d}\times\RR^{d})$. We can expand $K(x,y)$ out in
this basis as
\begin{equation}
K(x,y)=\sum^{\infty}_{k=1}\sum^{\infty}_{\ell=1}a_{k,\ell}\varphi_{k}(x)\varphi_{\ell}(y),
\end{equation}
with
\begin{equation}
\sum_{k,\ell}|a_{k,\ell}|^{2}<\infty
\end{equation}
since
\begin{equation}
\int|K|^{2}=\int|\sum_{k,\ell}a_{k,\ell}\varphi_{k}\varphi_{\ell}|^{2}=\sum_{k,\ell}|a_{k,\ell}|,
\end{equation}
and this must be finite. We just need to show $T$ is a limit of
finite-rank operators. Define
\begin{equation}
K_{n}(x,y):=\sum^{n}_{k=1}\sum^{n}_{\ell=1}a_{k,\ell}\varphi_{k}(x)\varphi_{\ell}(y),
\end{equation}
and we find
\begin{equation}
\|K-K_{n}\|_{L^{2}}^{2}=\sum_{\substack{k\geq n\\\text{or }\ell\geq n}}|a_{k,\ell}|^{2}\xrightarrow{n\to\infty}0.
\end{equation}
So then using Equation~\eqref{eq:math-108b:winter2026:lecture25:star},
\begin{equation}
\|T-T_{n}\|\leq\|K-K_{n}\|\xrightarrow{n\to\infty}0,
\end{equation}
where
\begin{equation}
(T_{n}f)(x) = \int_{\RR^{d}}K_{n}(x,y)f(y)\,\D y
\end{equation}
is a finite-rank operator. Hence $T$ is the norm limit of a sequence
of finite-rank operators and $T$ is bounded. This means $T$ is compact.
\end{proof}

\begin{remark}
More generally, if we are working on a locally compact Hausdorff space
$X$ equipped with a positive Borel measure, then we can define a
notion of a Hilbert--Schmidt integral operator. In this case, if
$L^{2}(X)$ is separable, and the kernel $K\in L^{2}(X\times X)$, then
the operator $T\colon L^{2}(X)\to L^{2}(X)$ defined by
\begin{equation}
(Tf)(x)=\int_{X}K(x,y)f(y)\,\D y
\end{equation}
is compact.
\end{remark}

\begin{definition}
We can define the \define{Hilbert--Schmidt Norm} for an integral
operator $T$ on a Hilbert space $\mathcal{H}=L^{2}(\RR^{d})$ as
\begin{equation}
\|T\|_{\text{HS}}:=\|K\|_{L^{2}(\RR^{d}\times\RR^{d})}.
\end{equation}
\end{definition}

\begin{remark}
More generally, if $T$ is any linear operator on a Hilbert space
$\mathcal{H}$, then we can define its Hilbert--Schmidt norm to be
\begin{equation}
\|T\|_{\text{HS}}^{2}:=\sum_{i}\|Te_{i}\|_{\mathcal{H}}^{2}
\end{equation}
where $\{e_{i}\}_{i\in I}$ is an orthonormal basis. This need not be a
countable basis, we only require the sum contains at most countably
many non-zero terms. This is independent of choice of orthonormal basis.

Also, on a finite-dimensional Hilbert space, this just coincides with
the Frobenius norm.
\end{remark}

\begin{fact}
We have $\|T\|_{\text{op}}\leq\|T\|_{\text{HS}}$.
\end{fact}

\subsection{Banach Spaces}

\begin{definition}
A \define{Banach Space} is a complete normed vector space, where
``complete'' as a metric space with the metric is the induced one $d(x,y):=\|x-y\|$.
\end{definition}

\begin{node}
We have two aspects to Banach spaces: they are metric spaces, and they
are measure spaces.

Metric spaces had nowhere dense (closure has empty interior) and
meager sets (countable unions of nowhere dense sets) as its ``small sets''.
Residual (= co-meager) sets were ``large'' for metric spaces.

Measure spaces had null sets be its ``small'' sets. And full measure
(= co-null) sets were ``large''.

How do these two notions of ``small'' sets interact with Banach
spaces? Are all meager sets also null? Are all null sets also meager?
\end{node}

\begin{observation}
Both null and meager sets are closed under countable unions. Is it too
much to hope that they ``measure'' the ``same'' kind of smallness?
\end{observation}

\begin{example}[Open dense trap]
Let $\{q_{n}\}_{n\in\NN}$ be an enumeration of $\QQ$. Fix
$\varepsilon>0$. Let
\begin{equation}
U_{\varepsilon}=\bigcup_{n=1}^{\infty}(q_{n}-\varepsilon 2^{-n},q_{n}+\varepsilon2^{-n})
\end{equation}
be a union of countably many open intervals, each interval has size
$2\cdot\varepsilon2^{-n}=\varepsilon2^{-n+1}$. We know
$U_{\varepsilon}$ is dense in $\RR$ since $\QQ\subset U_{\varepsilon}$.
We also know $U_{\varepsilon}$ is an open set. Its complement
$U_{\varepsilon}^{\complement}$ is closed and has empty interior, so
it is nowhere dense.

Calculating the Lebegue measure $m(U_{\varepsilon})$ is part of
homework 4.

The moral; these two notions of smallness \underline{\emph{do not}} coincide.
\end{example}

\subsection*{Back to Banach spaces in general\dots}

\begin{fact}
In an infinite-dimensional Banach space, we cannot define a
translation-invariant measure on them. But their topology
survives. The baire category theorem governs those meager and residual sets.
\end{fact}

\begin{theorem}[Uniform boundedness principle]
Let $X$ be a Banach space, let $Y$ be a normed vector space, let
$\{T_{\alpha}\colon X\to Y\}_{\alpha\in A}$ be a family of continuous
linear operators.

If we have pointwise boundedness ($\forall x\in X\ldotp\sup_{\alpha\in A}\|T_{\alpha}x\|_{Y}<\infty$),
then we have \define{Uniform Boundedness}, i.e., $\sup_{\alpha\in A}\|T_{\alpha}\|_{\text{op}}<\infty$.
\end{theorem}

Note the converse is true since
$\|T_{\alpha}x\|_{Y}\leq\|T_{\alpha}\|_{\text{op}}\|x\|_{X}$. 